\scanpage{16}

\begin{center}
U N I V E R S I T E\\
DES SCIENCES ET TECHNIQUES\\
DU LANGUEDOC\\
\textasciicircum\textasciicircum\textasciicircum\textasciicircum\textasciicircum\textasciicircum\textasciicircum\textasciicircum\textasciicircum\textasciicircum\textasciicircum\textasciicircum\textasciicircum
\end{center}

Montpellier, le 14 octobre 1983

Service de la Scolarité

Diffusion : Liste A\\
MA/gs \hfill N° 4469

\begin{center}
AVIS DE SOUTENANCE D'UNE THESE DE DOCTORAT DE 3e CYCLE

\uuline{`` SCIENCES AGRONOMIQUES option PHYTOTECHNIE ''}
\end{center}

\uline{NOM et Prénom du candidat} : LECONTE Antoine

\uline{Titre de la Thèse} : ``LA REPRODUCTION SEXUEE DE L'HEVEA BRASILIENSIS (UNE APPROCHE HISTOLOGIQUE ET EXPERIMENTALE)''

\uline{Date, heure et lieu de soutenance prévus} : 27 octobre 1983 à 9h.30 - Salle de cours n°9 - (Bât. administratif ) Université des Sciences et Techniques du Languedoc.

\uline{Nom du Directeur de thèse} : MM. Charles HEBANT - BONNET

\uline{Désignation du laboratoire où la thèse est préparée} : Laboratoire de Biologie végétale - U.S.T.L.

\uline{Analyse de l'ouvrage} : L'amélioration génétique de l'Hévéa est entravée par un très mauvais taux de réussite des pollinisations contrôlées.

L'étude histologique montre que le processus reproducteur (gaméto-génèse, fécondation et développement embryonnaire) se déroule normalement. Une exception est signalée : le cas du clone GT 1 dont la stérilité mâle est décrite. La fécondation a lieu moins de 24 heures après la pollinisation.

L'étude expérimentale, menée en Côte d'Ivoire, porte sur les modalités de la floraison et de la fructification en conditions naturelles et artificielles. La majorité des fleurs subissent une élimination naturelle dans un délai de 5 à 15 jours après la pollinisation. Des traitements hormonaux (ANA, chlorure de cobalt) peuvent augmenter significativement le taux de fructification.

Une explication physiologique de la faible fertilité de l'Hévéa est proposée.

\uline{Durée approximative des recherches ayant abouti à la thèse} : 2 ans

\uline{Composition du jury proposé} : \begin{tabular}[t]{ll}
MM. & BONNET\\
& D'AUZAC\\
& HALLE\\
& NICOLAS D.
\end{tabular}

\begin{flushright}
P/ le Secrétaire Général,\\
l'Attachée Principale d'Administration scolaire et universitaire,

\unsure{A. Gentil}

\uline{A. GENTIL}
\end{flushright}

\margin{15}
